% © 2025 Ing. David Jaroš — CC BY-NC-ND 4.0
%
% This work is licensed under a Creative Commons Attribution-NonCommercial-NoDerivatives
% 4.0 International License (CC BY-NC-ND 4.0).
%
% License History: Earlier drafts (up to v0.3) were released under CC BY 4.0.
% From v0.4 onward, all material is released under CC BY-NC-ND 4.0 to protect
% the integrity of the theoretical work during ongoing academic development.
%
% See LICENSE.md for full license text.

% UBT-AI-PROVENANCE-BEGIN
% schema: ubt-ai-provenance/v1
% tier: B_machine_verified
% ai_assistance: disclosed
% human_review: machine-verification
% editorial_responsibility: Ing. David Jaroš
% policy: ../../AI_PROVENANCE.md
% notice: Machine-verified against named sources or verifiers; individual attestation is not claimed.
% UBT-AI-PROVENANCE-END

\ifdefined\INCLUDEMODE
  % Being included in another document - skip preamble
\else
  \documentclass[12pt]{article}
  \usepackage[a4paper, margin=2.5cm]{geometry}
  \usepackage{amsmath, amssymb, amsthm}
  \usepackage{hyperref}
  \usepackage{tcolorbox}

  \newtheorem{theorem}{Theorem}[section]
  \newtheorem{lemma}[theorem]{Lemma}
  \newtheorem{corollary}[theorem]{Corollary}
  \newtheorem{proposition}[theorem]{Proposition}
  \theoremstyle{definition}
  \newtheorem{definition}[theorem]{Definition}
  \theoremstyle{remark}
  \newtheorem{remark}[theorem]{Remark}

  \title{\textbf{GR Closure Step 3: Lorentzian Signature Theorem\\
         Complex-Time Direction with Lorentzian Projection Convention}}
  \author{Ing. David Jaroš}
  \date{May 2026}

  \begin{document}
  \maketitle
\fi

\section{Lorentzian Signature from Complex-Time Direction and Lorentzian Projection}
\label{sec:signature_theorem}

\begin{tcolorbox}[colback=blue!5!white,colframe=blue!75!black,title=Scope]
This step proves Lorentzian signature in the \emph{projected GR sector}.  
The result is not a claim that complex time alone fixes signature for arbitrary field data.
\end{tcolorbox}

\subsection{Algebraic carrier and projection setting}

The biquaternion algebra is
\begin{equation}
\mathbb B = \mathbb C\otimes_{\mathbb R}\mathbb H \cong \mathrm{Mat}(2,\mathbb C),
\qquad \dim_{\mathbb R}\mathbb B = 8.
\end{equation}

\begin{remark}[Clifford relation statement]
$\mathbb B$ provides the spinorial / even-Clifford algebraic carrier naturally associated with the Lorentzian sector.
It is not the full real Clifford algebra $\mathrm{Cl}_{1,3}(\mathbb R)$, whose real dimension is $16$.
\end{remark}

For the projected GR sector, define the Lorentzian bilinear form by
\begin{equation}
\label{eq:eta-bilinear-step3}
\langle A,B\rangle_\eta := \operatorname{Re}\operatorname{Tr}\!\left(A^\dagger\,\eta_{\mathcal B}\,B\right),
\end{equation}
where $\eta_{\mathcal B}$ is an external Lorentzian bilinear-form operator on the selected real four-dimensional projection subspace (equivalently a representation-level metric operator), not an element of $\mathbb B$.

\subsection{Clifford algebra, Pauli representation, and signature convention}

The Clifford relation is
\begin{equation}
\{\gamma^\mu,\gamma^\nu\}=2\eta^{\mu\nu}I.
\end{equation}
For signature $(-,+,+,+)$,
\begin{equation}
(\gamma^0)^2=-I,
\qquad
(\gamma^i)^2=+I.
\end{equation}
Pauli matrices satisfy
\begin{equation}
\sigma_i^2=I,
\end{equation}
and quaternionic imaginary units are represented by
\begin{equation}
i\sigma_i,
\qquad
(i\sigma_i)^2=-I.
\end{equation}

\begin{remark}[Important warning]
The set $\{iI,\sigma_1,\sigma_2,\sigma_3\}$ has the desired squares for $(-,+,+,+)$,
but it does not by itself form a full Clifford basis because $iI$ commutes with $\sigma_i$.
Correct squares alone are not sufficient; Clifford anticommutation must also hold.
\end{remark}

\subsection{Lorentzian admissibility}

For $\Theta: M^4\times\mathbb C_\tau\to\mathbb B$, define
\begin{equation}
M_{\mu\nu}(\Theta):=\langle \partial_\mu\Theta,\partial_\nu\Theta\rangle_\eta.
\end{equation}
Lorentzian admissibility requires:
\begin{enumerate}
  \item real linear independence of $\{\partial_\mu\Theta\}_{\mu=0}^3$ on the selected projection sector,
  \item non-vanishing Lorentzian Gram determinant
  \begin{equation}
  \det\!\big(M_{\mu\nu}(\Theta)\big)\neq 0.
  \end{equation}
\end{enumerate}

\subsection{Main theorem}

\begin{theorem}[Lorentzian Signature from Complex-Time Direction and Lorentzian Projection]
\label{thm:axiomB_signature}
\label{thm:signature}
Given AXIOM-B and the specified Lorentzian projection convention on the selected real four-dimensional sector,
the induced metric
\begin{equation}
\label{eq:induced-metric-step3}
g_{\mu\nu}=\frac{\langle \partial_\mu\Theta,\partial_\nu\Theta\rangle_\eta}{\mathcal N},
\qquad
\mathcal N=\left|\langle \partial_0\Theta,\partial_0\Theta\rangle_\eta\right|>0,
\end{equation}
has signature $(-,+,+,+)$ for Lorentzian-admissible $\Theta$.
\end{theorem}

\begin{proof}
AXIOM-B selects the timelike complex-time direction. Under the Lorentzian projection convention,
$\langle \partial_0\Theta,\partial_0\Theta\rangle_\eta<0$ on the admissible branch, while projected spatial directions have positive norm.
By Lorentzian admissibility, $\det(M_{\mu\nu})\neq0$, so the Gram matrix is non-degenerate in the indefinite setting.
Positive normalization by $\mathcal N$ rescales but does not change inertia indices, hence the signature is $(-,+,+,+)$.
\end{proof}

\begin{remark}[Normalization role]
$\mathcal N$ is scale-fixing only. It does not determine the signature.
\end{remark}

\begin{remark}[Claim boundary]
This theorem is a structural result within the chosen Lorentzian projection convention.
It should be stated as: signature from AXIOM-B plus Lorentzian projection/admissibility,
not AXIOM-B alone.
\end{remark}

\ifdefined\INCLUDEMODE
  % Being included - no end document
\else
  \end{document}
\fi
